% ================================================================
%  TAREAS — Semillero de Astronomía
%  Institución Educativa Enrique Olaya Herrera — Medellín
% ================================================================
\documentclass[12pt,oneside]{book}

% -------------------------------------------------
% PAQUETES
% -------------------------------------------------
\usepackage[spanish,es-nodecimaldot]{babel}
\usepackage[utf8]{inputenc}
\usepackage[T1]{fontenc}
\usepackage{lmodern}
\usepackage[a4paper,margin=2.5cm]{geometry}
\usepackage{microtype}
\usepackage{setspace}
\usepackage{graphicx}
\usepackage{xcolor}
\usepackage{fancyhdr}
\usepackage{titlesec}
\usepackage{hyperref}
\usepackage{bookmark}
\usepackage{tocloft}
\usepackage{tcolorbox}
\tcbuselibrary{theorems, skins, breakable}
\usepackage{amsmath,amssymb}
\usepackage{enumitem}
\usepackage{csquotes}

% -------------------------------------------------
% CONFIGURACIÓN VISUAL  (paleta cósmica)
% -------------------------------------------------
\definecolor{azulCosmos}{HTML}{1B3A6B}
\definecolor{azulCielo}{HTML}{2980B9}
\definecolor{doradoEstrella}{HTML}{D4A017}
\definecolor{grisClaro}{HTML}{F4F6F8}
\definecolor{grisTexto}{HTML}{2E2E2E}
\definecolor{grisSec}{HTML}{6C757D}
\definecolor{colorConcepto}{HTML}{EBF5FB}
\definecolor{colorActividad}{HTML}{EAFAF1}
\definecolor{colorIdea}{HTML}{FEF9E7}
\definecolor{colorNota}{HTML}{F2F3F4}
\definecolor{colorCompromiso}{HTML}{EAF0FB}

\hypersetup{
    colorlinks=true,
    linkcolor=azulCosmos,
    urlcolor=azulCielo,
    citecolor=azulCosmos,
    pdftitle={Tareas — Semillero de Astronomía},
    pdfauthor={Daniel Soto Villada}
}

\setstretch{1.12}
\setlength{\parskip}{0.4em}
\setlength{\parindent}{0pt}
\setcounter{secnumdepth}{3}
\setcounter{tocdepth}{2}
\setlength{\headheight}{26pt}

% -------------------------------------------------
% ESTILOS DE TÍTULOS
% -------------------------------------------------
\titleformat{\chapter}[display]
  {\bfseries\Huge\color{azulCosmos}}
  {\filleft\Large\chaptername\ \thechapter}
  {1ex}
  {\titlerule\vspace{1ex}\filright}
  [\vspace{1ex}\titlerule]

\titleformat{\section}
  {\bfseries\Large\color{azulCosmos}}
  {\thesection}{0.7em}{}

\titleformat{\subsection}
  {\bfseries\large\color{grisTexto}}
  {\thesubsection}{0.6em}{}

\titleformat{\subsubsection}
  {\bfseries\normalsize\color{grisTexto}}
  {\thesubsubsection}{0.5em}{}

% -------------------------------------------------
% ENCABEZADOS Y PIES
% -------------------------------------------------
\pagestyle{fancy}
\fancyhf{}
\fancyhead[L]{\textcolor{grisSec}{\nouppercase{\leftmark}}}
\fancyhead[R]{\textcolor{grisSec}{Semillero de Astronomía — IEEOH}}
\fancyfoot[C]{\textcolor{grisSec}{\thepage}}
\renewcommand{\headrulewidth}{0.4pt}
\renewcommand{\headrule}{%
  \hbox to\headwidth{\color{azulCosmos}\leaders\hrule height \headrulewidth\hfill}}

% -------------------------------------------------
% ÍNDICE
% -------------------------------------------------
\renewcommand{\cftchapfont}{\bfseries\color{azulCosmos}}
\renewcommand{\cftsecfont}{\color{grisTexto}}
\renewcommand{\cftchappagefont}{\bfseries\color{azulCosmos}}
\renewcommand{\cftsecpagefont}{\color{grisTexto}}

% -------------------------------------------------
% ENTORNOS TCOLORBOX
% -------------------------------------------------

% Caja genérica destacada
\newtcolorbox{destacado}{
  colback=grisClaro,
  colframe=azulCosmos,
  boxrule=0.8pt,
  arc=2mm,
  left=3mm, right=3mm, top=1.5mm, bottom=1.5mm
}

% Concepto astronómico
\newtcbtheorem[number within=chapter]{concepto}{Concepto}{
  enhanced,
  colback=colorConcepto,
  colframe=azulCosmos,
  coltitle=white,
  colbacktitle=azulCosmos,
  attach boxed title to top left={yshift=-2mm, xshift=4mm},
  boxed title style={enhanced, arc=1mm, boxrule=0pt, colframe=azulCosmos},
  fonttitle=\bfseries,
  breakable,
  arc=1.5mm,
  boxrule=0.8pt,
  left=4mm, right=4mm, top=5mm, bottom=2mm,
  separator sign none,
  description delimiters={(}{)}
}{con}

% Actividad
\newtcbtheorem[number within=chapter]{actividad}{Actividad}{
  enhanced,
  colback=colorActividad,
  colframe=azulCosmos,
  coltitle=white,
  colbacktitle=azulCosmos,
  attach boxed title to top left={yshift=-2mm, xshift=4mm},
  boxed title style={enhanced, arc=1mm, boxrule=0pt, colframe=azulCosmos},
  fonttitle=\bfseries,
  breakable,
  arc=1.5mm,
  boxrule=0.8pt,
  left=4mm, right=4mm, top=5mm, bottom=2mm,
  separator sign none,
  description delimiters={(}{)}
}{act}

% Problema
\newtcbtheorem[number within=chapter]{problema}{Problema}{
  enhanced,
  colback=white,
  colframe=azulCosmos,
  coltitle=white,
  colbacktitle=azulCosmos,
  attach boxed title to top left={yshift=-2mm, xshift=4mm},
  boxed title style={enhanced, arc=1mm, boxrule=0pt, colframe=azulCosmos},
  fonttitle=\bfseries,
  breakable,
  arc=1.5mm,
  boxrule=0.8pt,
  left=4mm, right=4mm, top=5mm, bottom=2mm,
  separator sign none,
  description delimiters={(}{)}
}{prob}

% Solución paso a paso
\newtcbtheorem[number within=chapter]{solucion}{Solución}{
  enhanced,
  colback=colorActividad,
  colframe=azulCielo,
  coltitle=white,
  colbacktitle=azulCielo,
  attach boxed title to top left={yshift=-2mm, xshift=4mm},
  boxed title style={enhanced, arc=1mm, boxrule=0pt, colframe=azulCielo},
  fonttitle=\bfseries,
  breakable,
  arc=1.5mm,
  boxrule=0.8pt,
  left=4mm, right=4mm, top=5mm, bottom=2mm,
  separator sign none,
  description delimiters={(}{)}
}{sol}

% Idea / Propuesta
\newtcbtheorem[number within=chapter]{idea}{Idea}{
  enhanced,
  colback=colorIdea,
  colframe=doradoEstrella,
  coltitle=white,
  colbacktitle=doradoEstrella,
  attach boxed title to top left={yshift=-2mm, xshift=4mm},
  boxed title style={enhanced, arc=1mm, boxrule=0pt, colframe=doradoEstrella},
  fonttitle=\bfseries,
  breakable,
  arc=1.5mm,
  boxrule=0.8pt,
  left=4mm, right=4mm, top=5mm, bottom=2mm,
  separator sign none,
  description delimiters={(}{)}
}{idea}

% Nota / Observación
\newtcbtheorem[number within=chapter]{nota}{Nota}{
  enhanced,
  colback=colorNota,
  colframe=grisSec,
  coltitle=white,
  colbacktitle=grisSec,
  attach boxed title to top left={yshift=-2mm, xshift=4mm},
  boxed title style={enhanced, arc=1mm, boxrule=0pt, colframe=grisSec},
  fonttitle=\bfseries,
  breakable,
  arc=1.5mm,
  boxrule=0.6pt,
  left=4mm, right=4mm, top=5mm, bottom=2mm,
  separator sign none,
  description delimiters={(}{)}
}{nota}

% Compromiso personal
\newtcbtheorem[number within=chapter]{compromiso}{Compromiso}{
  enhanced,
  colback=colorCompromiso,
  colframe=azulCielo,
  coltitle=white,
  colbacktitle=azulCielo,
  attach boxed title to top left={yshift=-2mm, xshift=4mm},
  boxed title style={enhanced, arc=1mm, boxrule=0pt, colframe=azulCielo},
  fonttitle=\bfseries,
  breakable,
  arc=1.5mm,
  boxrule=0.8pt,
  left=4mm, right=4mm, top=5mm, bottom=2mm,
  separator sign none,
  description delimiters={(}{)}
}{comp}

% -------------------------------------------------
% DATOS DEL DOCUMENTO
% -------------------------------------------------
\newcommand{\NombreSemillero}{Semillero de Astronomía}
\newcommand{\Institucion}{Institución Educativa Enrique Olaya Herrera}
\newcommand{\Ciudad}{Medellín, Colombia}
\newcommand{\AsesorPrincipal}{Daniel Soto Villada}
\newcommand{\CoAsesor}{Gerardo Ríos}
\newcommand{\TituloDocumento}{Tareas del\\[0.4em]Semillero de Astronomía}

% -------------------------------------------------
% DOCUMENTO
% -------------------------------------------------
\begin{document}
\hypersetup{pageanchor=false}

% ---------- PORTADA ----------
\begin{titlepage}
    \centering
    \vspace*{1.5cm}

    {\Large\textcolor{grisSec}{\Institucion}\par}
    \vspace{0.2cm}
    {\large\textcolor{grisSec}{\Ciudad}\par}

    \vspace{2.0cm}

    {\Huge\bfseries\textcolor{azulCosmos}{\TituloDocumento}\par}
    \vspace{0.5cm}
    {\Large\textcolor{grisSec}{Problemas y Ejercicios de Astronomía}\par}

    \vspace{2.4cm}

    \begin{tcolorbox}[
        colback=grisClaro,
        colframe=azulCosmos,
        boxrule=0.9pt,
        arc=2mm,
        width=0.82\textwidth
    ]
        \centering
        {\large\bfseries\textcolor{azulCosmos}{Asesores}\par}
        \vspace{0.3cm}
        {\large \AsesorPrincipal\ \textcolor{grisSec}{(Director Académico)}\par}
        \vspace{0.1cm}
        {\large \CoAsesor\ \textcolor{grisSec}{(Director de Gestión)}\par}
    \end{tcolorbox}

    \vspace{1.6cm}

    \begin{tcolorbox}[
        colback=colorIdea,
        colframe=doradoEstrella,
        boxrule=0.7pt,
        arc=2mm,
        width=0.82\textwidth
    ]
        \centering
        \itshape\large\textcolor{grisTexto}{%
        ``La astronomía compele al alma a mirar hacia arriba y nos lleva\\
        de este mundo a otro.''}\\[0.4em]
        \normalsize\textcolor{grisSec}{--- Platón}
    \end{tcolorbox}

    \vfill
    {\large\textcolor{grisSec}{\today}\par}
    \vspace{0.5cm}
\end{titlepage}

% ---------- PÁGINAS PRELIMINARES ----------
\frontmatter
\tableofcontents
\clearpage

% ---------- CONTENIDO ----------
\clearpage
\hypersetup{pageanchor=true}
\mainmatter
\chapter{Tarea 1 — Sistema Solar (resuelta)}

\begin{destacado}
Este documento fue reorganizado para una lectura más clara:
cada \textbf{Problema} está seguido inmediatamente por su
\textbf{Solución paso a paso}.
\end{destacado}

\section{Órbita del Sistema Solar y escala humana}

\begin{problema}{Fracción de órbita recorrida por la humanidad}{prob:orbita-humana}
Desde la aparición de los \textit{Homo sapiens} (aprox. $300\,000$ años):
\begin{itemize}[leftmargin=1.5em]
  \item el Sistema Solar tarda $230$ millones de años en completar una órbita galáctica;
  \item su velocidad orbital promedio es $828\,000\ \mathrm{km/h}$.
\end{itemize}
Calcular:
\begin{enumerate}[leftmargin=1.5em, label=\alph*)]
  \item fracción de órbita recorrida;
  \item distancia recorrida en kilómetros;
  \item ángulo barrido (en grados), asumiendo órbita circular.
\end{enumerate}
\end{problema}

\begin{solucion}{Paso a paso}{sol:orbita-humana}
\textbf{1) Fracción de órbita}
\[
f=\frac{300\,000}{230\,000\,000}=1.304\times 10^{-3}
\]
Esto equivale a:
\[
f\cdot 100=0.1304\%
\]

\textbf{2) Distancia recorrida}
\[
t=300\,000\ \text{años}\times 365.25\ \frac{\text{días}}{\text{año}}\times 24\ \frac{\text{h}}{\text{día}}
=2.6298\times 10^{9}\ \text{h}
\]
Aquí se usa \(365.25\) días por año para incluir años bisiestos de forma aproximada.
Es el promedio de \(365+0.25\) días, considerando un día adicional cada cuatro años.
\[
d=v\,t=(828\,000\ \mathrm{km/h})(2.6298\times10^9\ \mathrm{h})
\approx 2.18\times10^{15}\ \mathrm{km}
\]

\textbf{3) Ángulo barrido}
\[
\theta=f\cdot 360^\circ=(1.304\times10^{-3})(360^\circ)\approx 0.47^\circ
\]

\textbf{Resultado final:} en el tiempo de existencia humana moderna, el Sistema Solar
ha recorrido cerca de \textbf{0.13\%} de su órbita galáctica, una distancia aproximada
de \(\mathbf{2.18\times10^{15}}\) km, y un ángulo de \(\mathbf{0.47^\circ}\).
\end{solucion}

\section{Escalas de distancia en el Sistema Solar}

\begin{problema}{Escala y tiempos de viaje}{prob:escala-viaje}
Usar:
\[
1\ \mathrm{AU}\approx 150\times 10^6\ \mathrm{km}
\]
Esta es una aproximación estándar de la Unidad Astronómica para cálculos introductorios,
y la escala del modelo es:
\[
1\ \mathrm{AU}\longleftrightarrow 1\ \mathrm{cm}
\]
Calcular:
\begin{enumerate}[leftmargin=1.5em, label=\alph*)]
  \item distancia en la maqueta para Neptuno (\(30.1\ \mathrm{AU}\)),
        Heliosfera (\(80\) a \(100\ \mathrm{AU}\)),
        inicio de Nube de Oort (\(5\,000\ \mathrm{AU}\))
        y borde externo (\(100\,000\ \mathrm{AU}\));
  \item tiempo de viaje a \(900\ \mathrm{km/h}\) para llegar a
        la Luna (\(384\,400\ \mathrm{km}\)),
        Neptuno (\(\approx 30\ \mathrm{AU}\)),
        Heliosfera (\(100\ \mathrm{AU}\))
        y borde externo de la Nube de Oort (\(100\,000\ \mathrm{AU}\)).
\end{enumerate}
\end{problema}

\begin{solucion}{Paso a paso}{sol:escala-viaje}
\textbf{1) Distancias en la maqueta (\(1\ \mathrm{AU}=1\ \mathrm{cm}\))}
\begin{itemize}[leftmargin=1.5em]
  \item Neptuno: \(30.1\ \mathrm{AU}\Rightarrow 30.1\ \mathrm{cm}\).
  \item Heliosfera: \(80\) a \(100\ \mathrm{AU}\Rightarrow 80\) a \(100\ \mathrm{cm}\) (\(0.8\) a \(1\ \mathrm{m}\)).
  \item Inicio Nube de Oort: \(5\,000\ \mathrm{AU}\Rightarrow 5\,000\ \mathrm{cm}=50\ \mathrm{m}\).
  \item Borde externo Nube de Oort: \(100\,000\ \mathrm{AU}\Rightarrow 100\,000\ \mathrm{cm}=1\,000\ \mathrm{m}=1\ \mathrm{km}\).
\end{itemize}

\textbf{2) Tiempos de viaje a \(900\ \mathrm{km/h}\)}
\[
t=\frac{d}{v}
\]
\begin{itemize}[leftmargin=1.5em]
  \item Luna:
  \[
  t=\frac{384\,400}{900}=427.1\ \mathrm{h}\approx 17.8\ \text{días}
  \]
  \item Neptuno (\(\approx 30\ \mathrm{AU}\)):
  \[
  d\approx 30(150\times10^6)=4.5\times10^9\ \mathrm{km}
  \]
  \[
  t=\frac{4.5\times10^9}{900}=5.0\times10^6\ \mathrm{h}\approx 570\ \text{años}
  \]
  \item Heliosfera (\(100\ \mathrm{AU}\)):
  \[
  d=100(150\times10^6)=1.5\times10^{10}\ \mathrm{km}
  \]
  \[
  t=\frac{1.5\times10^{10}}{900}=1.67\times10^7\ \mathrm{h}\approx 1\,900\ \text{años}
  \]
  \item Borde externo Nube de Oort (\(100\,000\ \mathrm{AU}\)):
  \[
  d=100\,000(150\times10^6)=1.5\times10^{13}\ \mathrm{km}
  \]
  \[
  t=\frac{1.5\times10^{13}}{900}=1.67\times10^{10}\ \mathrm{h}\approx 1.9\times10^6\ \text{años}
  \]
\end{itemize}

\textbf{Resultado final:} la maqueta muestra una enorme diferencia de escalas
(de centímetros a kilómetros), y los tiempos de viaje con tecnología convencional
se vuelven rápidamente inabordables fuera del entorno planetario cercano.
\end{solucion}

% ---------- FUTURAS TAREAS (se irán agregando aquí) ----------
% \input{Tareas/tarea02}
% \input{Tareas/tarea03}
% ...

\end{document}
